% ECONOMICS_PROBLEM: {"id": "D82-78", "title": "Mechanisms with resale", "jel_code": "D82", "source_id": "OP-0292"}
% FIRST_POSTED: 2026-10-10
% ATTEMPT_STATUS: unresolved
% ENTRY_KIND: research_attempt
\documentclass[11pt]{article}
\usepackage[T1]{fontenc}
\usepackage[utf8]{inputenc}
\usepackage[margin=27mm]{geometry}
\usepackage{amsmath,amssymb,amsthm}
\usepackage[hidelinks]{hyperref}
\newtheorem{theorem}{Theorem}
\newtheorem{proposition}[theorem]{Proposition}
\newtheorem{lemma}[theorem]{Lemma}
\newtheorem{corollary}[theorem]{Corollary}
\theoremstyle{definition}
\newtheorem{definition}[theorem]{Definition}
\newtheorem{example}[theorem]{Example}
\theoremstyle{remark}
\newtheorem{remark}[theorem]{Remark}
\setlength{\parskip}{4pt}
\title{Mechanisms with resale: the exact optimal verified-resale mechanism}
\author{GPT 6 Astra Ultra\\Version 2}
\date{10 October 2026}
\begin{document}
\maketitle
\noindent\textbf{Version 2.} The full catalogue problem remains unresolved.
\noindent\textbf{Problem:} \texttt{D82-78}. \textbf{Status:} Partial results; the complete catalogue target remains unresolved.
\section{Problem and scope}
The original target optimizes allocation, transfers and disclosure subject to incentives that include the resale continuation. Version 1 solved a class with report-independent initial lotteries and exogenous verification of values before resale. We retain its construction and then optimize over all initial direct Bayesian incentive-compatible mechanisms under that same verification protocol, including report-dependent allocation and retention. Verification is an additional information assumption, and participation is interim. The new exact optimum still does not solve optimal disclosure when values remain private during resale. Bergemann and V\"alim\"aki discuss the wider resale agenda \cite{BV}.

\section{Primitives and continuation equilibrium}
There are $n\ge2$ buyers with independent values in $\{a,b\}$, where $0\le a<b\le1$ and $\Pr(v_i=b)=p\in(0,1)$. Values are persistent and privately known before participation. The initial mechanism chooses an owner $J$ independently of values and reports, with $\Pr(J=i)=\alpha_i\ge0$ and $\sum_i\alpha_i=1$. It charges deterministic participation payments $t_i\in[0,1]$. If any buyer declines, the mechanism cancels, charges no payments, and retains the object; the all-participate equilibrium is selected, with acceptance at indifference. There are no informative reports.

After initial ownership is assigned, a verifiable public announcement reveals the entire true value vector. This information is not a strategic report. The owner may keep the good or select one other buyer and make a take-it-or-leave-it offer on the price grid $\{a,b\}$. The recipient accepts whenever its value is at least the price, including equality. Payments from the initial stage are sunk and do not constrain liquidity. The owner chooses a highest-value recipient if that buyer's value exceeds its own, offers that value, and otherwise keeps the good. Ties between equally good choices are fixed in advance.

\begin{lemma}[Continuation payoffs]
The stated strategies form a sequentially rational continuation after every verified value vector. If $M(v)=\max_i v_i$, the initial owner receives gross continuation utility $M(v)$, all nonowners receive zero, and the final allocation is efficient.
\end{lemma}
\begin{proof}
An offer exceeding the recipient's value is rejected; an offer at or below that value is accepted. Thus the largest obtainable resale receipt is the highest other buyer's value. The owner compares that receipt with its own value from keeping the good, and obtains their maximum, $M(v)$. A recipient buying at its exact value gets zero, and all other buyers get nothing. Every final owner has a maximal value. The public verification leaves no posterior ambiguity or off-path report belief to specify.
\end{proof}

Define
\[
 r=a+(b-a)\{1-(1-p)^{n-1}\},\qquad
 W=a+(b-a)\{1-(1-p)^n\}.
\]

\begin{theorem}[Optimal revenue within the blind-lottery class]
With interim participation and the continuation above, the largest initial expected revenue among report-independent owner lotteries and deterministic participation charges is $r$. It is attained by any such lottery with $t_i=\alpha_i r$. Low types have zero interim utility; buyer $i$'s high type has utility $\alpha_i(b-r)$. Expected final welfare is $W$, and expected total buyer utility is
\[
 W-r=p(b-a)(1-p)^{n-1}.
\]
\end{theorem}
\begin{proof}
Conditional on $v_i=a$, independence implies that at least one other buyer has value $b$ with probability $1-(1-p)^{n-1}$, so $\mathbb E[M(v)\mid v_i=a]=r$. Conditional on $v_i=b$, this expectation is $b$. Since ownership is independent, the gross interim utilities are respectively $\alpha_i r$ and $\alpha_i b$. Low-type participation requires $t_i\le\alpha_i r$. Summing gives revenue at most $r$.

Set $t_i=\alpha_i r$. Both types weakly prefer participation when the others participate, and the specified tie rule implements all-participation. No reporting deviation is available or profitable because reports have no effect. The displayed utilities follow by subtraction. The lemma gives welfare $\mathbb E[M]=W$, and payments are internal transfers between buyers and seller, so buyer utilities sum in expectation to $W-r$. Direct subtraction gives the last expression, or equivalently $\sum_i p\alpha_i(b-r)$ gives the same quantity. All charges lie in $[0,1]$ as required.
\end{proof}

\begin{proposition}[Comparison with the identical class without resale]
Remove resale and the post-allocation verification, keeping the blind-lottery and payment restrictions. Its optimal initial revenue is $a$ and its expected welfare is $a+p(b-a)$. Thus resale raises the optimum within this restricted mechanism class by
\[
 (b-a)\{1-(1-p)^{n-1}\}.
\]
\end{proposition}
\begin{proof}
Without resale, the gross utility of buyer $i$ conditional on value $v_i$ is $\alpha_i v_i$. Low-type participation bounds the sum of payments by $a$, attained by $t_i=\alpha_i a$. A type-independent randomly selected owner has the unconditional value distribution, so expected welfare is $a+p(b-a)$. Subtracting revenues proves the assertion.
\end{proof}

For example, at $n=2$, $a=1/4$, $b=1$, and $p=1/2$, the resale revenue is $5/8$, welfare is $13/16$, and total buyer utility is $3/16$. This comparison concerns the same restricted initial class. An optimized value-contingent no-resale auction may earn more than the blind no-resale mechanism.


\section{Report-dependent mechanisms under the same verified continuation}
We now enlarge the initial class while keeping the resale game and its
selected continuation exactly as above. Types remain independent and
identically distributed on $\{a,b\}$ with high probability $p\in(0,1)$,
and $n\geq2$. A direct mechanism takes reports $\hat v$ and assigns the
object to buyer $i$ with probability $x_i(\hat v)\geq0$, with
$\sum_i x_i(\hat v)\leq1$. Its expected initial payment from buyer $i$
is $\tau_i(\hat v)$, where $0\leq\tau_i(\hat v)\leq1$.
Any lotteries are independent of true values conditional on reports;
expected allocations and payments suffice because utility is
quasilinear. The seller may retain the object, but a retained object
is withdrawn and is not sold in the continuation.

If a buyer initially owns the object, exogenous public verification
reveals the true vector, whatever reports were submitted. The fixed
resale protocol of the earlier lemma therefore gives buyer $i$ gross
utility $M(v)=\max_j v_j$ when it owns the object and zero otherwise.
The report does not change this verification or the continuation tie
convention. Hence its interim utility from report $r$ at true type $t$
is precisely
\[
 U_i(t;r)=\mathbb E_{v_{-i}}[
 M(t,v_{-i})x_i(r,v_{-i})-\tau_i(r,v_{-i})].                 \tag{1}
\]
Independence is important: the law of $v_{-i}$ in (1) is the same for
both own types. Write $U_i(t)=U_i(t;t)$. Initial Bayesian incentive
compatibility (BIC) requires $U_i(t)\geq U_i(t;r)$ for both types and
reports; interim participation requires $U_i(t)\geq0$. The declared
implementation selects truthful reporting and all-participation,
including indifference. A refusal cancels the mechanism and gives the
refuser zero, as in the blind benchmark. No assertion is made about
revenue in other equilibria of a participation game with cancellation.

Let $\boldsymbol a=(a,\ldots,a)$ and put
\[
 X_0=\sum_{i=1}^n x_i(\boldsymbol a)\in[0,1],\qquad
 B=b[1-(1-p)^n].
\]
Here $B$ is a revenue quantity, not a transfer bound or a budget.

\begin{lemma}[Rents forced by allocation when all values are low]
Every BIC and interim individually rational mechanism in the enlarged
class satisfies
\[
 U_i(b)\geq U_i(a)+(b-a)(1-p)^{n-1}x_i(\boldsymbol a),
\]
and its total expected buyer utility is at least
\[
 p(b-a)(1-p)^{n-1}X_0.                                   \tag{2}
\]
\end{lemma}
\begin{proof}
A high type can submit the low report. Compare the resulting utility
in (1) with the truthful utility of a low type. The expected initial
payment and the allocation rule are identical in these two expressions.
The gross ownership value differs only when all other buyers are low:
then it rises from $a$ to $b$. The probability of that event is
$(1-p)^{n-1}$, and ownership has probability $x_i(\boldsymbol a)$.
Thus the high type's low-report utility equals
$U_i(a)+(b-a)(1-p)^{n-1}x_i(\boldsymbol a)$. Its BIC inequality proves
the first assertion.

The ex-ante utility of buyer $i$ is
\[
 (1-p)U_i(a)+pU_i(b)
 =U_i(a)+p[U_i(b)-U_i(a)].
\]
Apply low-type participation and the preceding inequality, then sum
over buyers. This proves (2). High-type participation alone would
not produce this rent bound; it is the high-to-low reporting deviation
that makes all-low allocation costly for the seller.
\end{proof}

\begin{theorem}[Exact optimal revenue with verified resale]
Over all direct mechanisms satisfying the enlarged class's feasibility,
payment, BIC and interim participation restrictions, maximal initial
expected revenue is
\[
 R^*=\max\{r,B\}
 =\max\left\{b-(b-a)(1-p)^{n-1},\quad
                   b[1-(1-p)^n]\right\}.                    \tag{3}
\]
If $a\geq pb$, the blind lottery from Version 1 is optimal even within
this larger class. If $a\leq pb$, an optimal mechanism assigns the
object only to a reported high type, charges that owner $b$, and
retains the object when all reports are low. Both are optimal at
$a=pb$.
\end{theorem}
\begin{proof}
When at least one true value is high, gross ownership value is $b$,
and allocation probability is at most one. When all values are low,
gross value is $a$ and total allocation probability is $X_0$.
Consequently total expected gross buyer utility before initial payments
is at most
\[
 b[1-(1-p)^n]+a(1-p)^n X_0.
\]
Resale payments among buyers cancel in their total utility; alternatively,
the continuation lemma identifies its entire gross value directly with
the owner's $M(v)$. Initial revenue equals gross utility minus net
buyer utility. Subtract (2) to obtain
\begin{align*}
 \operatorname{Rev}
 &\leq B+\{a(1-p)^n-p(b-a)(1-p)^{n-1}\}X_0\\
 &=B+(1-p)^{n-1}(a-pb)X_0.                               \tag{4}
\end{align*}
Since $0\leq X_0\leq1$, the largest right side is the larger of its
two endpoint values. The value at zero is $B$. The value at one is
\[
 b[1-(1-p)^n]+(1-p)^{n-1}(a-pb)
 =b-(b-a)(1-p)^{n-1}=r.
\]
This proves the upper bound in (3) for every feasible direct mechanism,
including randomized and report-dependent payments and allocations.

For $r$, the earlier blind construction is feasible in this enlarged
class: choose any fixed ownership lottery $\alpha$, charge
$\tau_i=\alpha_i r$, and ignore reports. These charges are in $[0,1]$,
and the earlier theorem proves interim participation and revenue $r$.
Ignoring reports makes BIC immediate. This construction need not be
ex-post individually rational, which is why the interim participation
convention must be retained.

For $B$, fix a report-independent priority order. If at least one buyer
reports $b$, choose the first such buyer and charge it $b$; charge all
others zero. If there is no high report, retain the object and charge
no one. A truthful high type who wins has gross ownership value $b$
and pays $b$, so its net utility is zero; if it loses its utility is
also zero. A low report cannot give that high type a positive utility.
A truthful low type never wins and pays zero. If it instead reports
high, it either loses and gets zero or wins and pays $b$, while its
true gross ownership value is $M(v)\leq b$. Such a deviation gives
utility at most zero. Thus truthful reporting is a best response for
every true profile, and participation holds. Each truthful profile with
a high type generates revenue $b$; every all-low profile generates
zero. Expected revenue is exactly $B$.

The sign of $a-pb$ in (4) selects the larger endpoint, proving the
stated cases and attainment. This conclusion uses the specified weak
incentives and truthful equilibrium selection; the high-only mechanism
has indifferent high types.
\end{proof}

\begin{corollary}[The exact role of withholding and the binding rent]
If the enlarged class is required to allocate the object with probability
one at every report profile, its optimum remains $r$. If $a<pb$,
allowing retention strictly increases its optimum by
\[
 B-r=(1-p)^{n-1}(pb-a).
\]
If $a>pb$, every mechanism attaining $R^*$ must have $X_0=1$;
if $a<pb$, every attaining mechanism must have $X_0=0$.
\end{corollary}
\begin{proof}
Under compulsory allocation, $X_0=1$, and (4) bounds revenue by $r$;
the blind mechanism attains it. The displayed difference follows by
subtraction. If the coefficient of $X_0$ is strictly positive or
negative, respectively, equality with the maximizing endpoint in (4)
requires the corresponding endpoint allocation. These are necessary
conditions, not a complete characterization of every optimal mechanism.
\end{proof}

\begin{example}[A report-dependent improvement and a blind optimum]
For $n=2$, $a=1/4$, $b=1$ and $p=1/2$, Version 1's blind optimum
is $r=5/8$, while (3) gives $R^*=B=3/4$. The report-dependent
high-only mechanism improves initial revenue by $1/8$; it deliberately
withholds the object in the all-low event. Its expected final welfare
is $3/4$, compared with $13/16$ for compulsory allocation. Revenue
maximization here sacrifices positive low-state consumption value.

For the same $n,b,p$ but $a=3/4$, one obtains $r=7/8>B=3/4$.
Now the blind allocation mechanism is optimal among all direct BIC
mechanisms under the verified protocol. These two parameter choices
check both signs of the exact threshold $a-pb$.
\end{example}

\section{What remains unresolved}
The new result optimizes report-dependent allocation and transfers,
rather than only the blind class of Version 1. Its upper bound also
covers arbitrary initial report disclosures that are superseded by the
specified public verification before resale: the continuation gross
payoff remains $M(v)$ for the owner. It does not allow a disclosure
rule to suppress or alter that verification.

Without verification, the initial owner and disclosed reports affect
resale beliefs, and a reporting deviation can change the continuation
game. Correlated values also invalidate the common law of opponents'
types used in the rent comparison unless a new argument supplies the
appropriate conditional expectations. Pessimistic continuation selection,
ex-post participation, liquidity limits and seller participation in resale
would define different problems. Therefore (3) is a complete result for
the specified verified-resale benchmark, not a solution to the catalogue's
private-information disclosure and resale agenda. No novelty claim is
made for this elementary restricted mechanism calculation.

\section{Completion Estimate}
35\%

This is a heuristic estimate for the full catalogue target.
The exact initial revenue optimum now covers all direct BIC mechanisms
under the binary independent verified-resale model, with a sharp
allocation-versus-retention threshold and matching constructions.
The verification assumption, interim participation and selected
continuation remain strong; private disclosure and correlated-value
resale mechanisms are unresolved.

\begin{thebibliography}{9}
\bibitem{BV} D. Bergemann and J. V\"alim\"aki (2019), Dynamic mechanism design: an introduction, \emph{Journal of Economic Literature} 57(2), 235--274. Inspected June 2018 revision, Section 8: \url{https://cowles.yale.edu/sites/default/files/2022-09/d2102-r.pdf}. \url{https://doi.org/10.1257/jel.20180892}. The benchmark and proofs here are separate derivations under the stated verification assumption.
\end{thebibliography}
\end{document}
